\chapter{Deferred Shading}\label{ch:deferred-shading}

\section{Overview}

When PBR data is present, the viewer shades in two phases: the volume and
primitive renderers write a \emph{G-buffer} (multiple render targets), and a
fullscreen pass (\code{deferred_shade_frag.wgsl},
\code{DeferredShadePipeline} in \code{crates/rmesh-postprocess/src/lib.rs})
combines it with analytic lights, deep-shadow-map transmittance, ambient
occlusion, and screen-space GI/reflections, then tone-maps. Everything is
linear \code{Rgba16Float} end to end; sRGB encoding happens only in the
final \code{blit.wgsl} copy to the swapchain.

\section{G-buffer layout}

Both producers --- the volume interval renderer
(\code{interval_fragment.wgsl}, see Chapter~\ref{ch:interval-renderer}) and the
opaque primitive pass (\code{primitive_mrt.wgsl} in
\code{rmesh-compositor}) --- write the \emph{same} four color attachments,
plus the hardware depth buffer. All volume channels are premultiplied by
$\alpha$ so hardware blending composites them; primitives write $\alpha = 1$.

\begin{longtable}{@{}llll@{}}
\toprule
Slot & Binding & Format & Contents \\
\midrule
albedo & 1 & Rgba16Float & $\text{albedo} \cdot \alpha,\ \alpha$ \\
aux0 (PBR) & 2 & Rgba16Float & $r\alpha,\ f_0\alpha,\ m\alpha,\ \alpha$ \\
normals & 3 & Rgba8Unorm & $(\tfrac{n}{2}{+}\tfrac12)\alpha,\ \alpha$ \\
depth & 4 & Rgba16Float & $\alpha E[z],\ \text{occl}\cdot\alpha,\ \alpha E[z^2],\ \alpha$ \\
hw depth & 6 & Depth32Float & hardware $z$ \\
\bottomrule
\end{longtable}

Conventions worth knowing:
\begin{itemize}
\item \textbf{Metallic sits in \code{aux0.b}}, dielectric $f_0$ in
  \code{aux0.g} (glTF metallic-roughness order); primitives write 0 there,
  and the deferred pass falls back to $f_0 = 0.04$.
\item \textbf{Normals} are the only non-float16 target: bias-encoded
  \code{Rgba8Unorm}. Volume normals are the interpolated \emph{scalar-field
  gradient} (density gradient), not a surface normal; primitives compute a
  TBN-mapped mesh normal. The deferred pass un-premultiplies, un-biases, and
  \emph{negates} --- converting COLMAP (Y-down, Z-forward) source
  conventions to wgpu (Y-up, Z-back).
\item \textbf{Depth appears three ways} in the system: hardware
  \code{Depth32Float}, the volume's premultiplied expected-termination depth
  (with $E[z^2]$ for variance), and the shadow maps' normalized $[0,1]$
  moment depth. They must not be conflated.
\end{itemize}

\section{World-position reconstruction}

The fullscreen pass (a single vertex-index triangle, \code{draw(0..3)})
reconstructs the shaded point per pixel by taking the \emph{closer} of the
volume's expected termination depth (slot 3, normalized by its own $\alpha$)
and the linearized hardware depth, converting back to NDC and multiplying by
the inverse view-projection.

\section{Lighting model}

Lights are a read-only storage buffer of up to 16 \code{GpuLight} records
(\code{rmesh-util/src/shared.rs}): position, type (0 point, 1 spot, 2
directional), color, intensity, pre-normalized direction, and inner/outer
cone cosines. The viewer collects them each frame from point-light and
spot-light primitives and from active \emph{flare projectiles} (the flare
gun's glTF model carries an offset light).

Shading is a metallic-roughness GGX BRDF:
\begin{itemize}
\item Fresnel: Schlick, $F = F_0 + (1-F_0)(1-\cos\theta)^5$, with
  $F_0 = (1-m)\,f_0 + m\,\text{albedo}$.
\item Distribution: GGX/Trowbridge-Reitz with $a = r^2$
  (roughness clamped to $[0.08, 1]$),
  $D = \dfrac{a^2}{\pi\big((n \cdot h)^2 (a^2 - 1) + 1\big)^2}$.
\item Visibility: height-correlated Smith GGX.
\item Diffuse: $k_d = (1 - F_{\text{view}})(1 - m)$ times albedo, no
  $1/\pi$ (matching the parent project's parametric relighting branch).
\end{itemize}

Per light: attenuation is $1/d^2$ for point/spot (smooth cone falloff
between the cone cosines for spots) and $1$ for directional; the
contribution is scaled by $T \cdot \text{atten}$ where $T$ is the deep
shadow map transmittance (Chapter~\ref{ch:shadow-rendering}). Two optimizations
land per pixel: an early \code{continue} when $n \cdot l \le 0$ skips the
BRDF \emph{and} the cubemap shadow sample for back-facing lights, and
$k_d \cdot \text{albedo}$ is folded once per pixel.

\section{Ambient, occlusion, and screen-space effects}

The indirect term combines:
\begin{itemize}
\item \textbf{GTAO} (\code{gtao.wgsl}, R8Unorm + temporal accumulation and
  bilateral blur), multiplied with the primitive occlusion channel and
  blended by an \code{ao_strength} uniform;
\item \textbf{ambient specular} via the Lazarov/Karis analytic
  environment-BRDF fit (\code{env_brdf_approx}) --- no BRDF LUT texture;
\item \textbf{SSGI} (\code{ssgi_compute.wgsl}) --- hemisphere rays marched
  against a Hi-Z min-pyramid (\code{hiz_linearize.wgsl} /
  \code{hiz_downsample.wgsl}, R32Float, fusing hardware and volume depth),
  sampling the \emph{previous frame's linear, pre-tonemap lit output} for
  one bounce. Misses (off-screen or step-budget exhausted) return
  \textbf{black}: SSGI is a pure on-screen bounce term, and the
  environment contribution lives in the composite instead;
\item \textbf{SSR} (\code{ssr_compute.wgsl}) --- a mirror ray for smooth
  surfaces; the indirect specular blends SSR toward the hemispherical
  \emph{environment} color by smoothness$^2$, and an SSR miss falls back
  to the environment sampled along the world reflection direction.
\end{itemize}

The composite treats visibility asymmetrically, on purpose: the SSGI
bounce already carries screen-space visibility in its rays, so it is
\emph{not} multiplied by AO (that would double-count occlusion); the
environment hemisphere term \emph{is} AO-occluded. Per pixel:
\begin{align*}
\text{env\_diffuse} &= \operatorname{mix}(\text{ground},\ \text{sky},\
  \tfrac{n_y+1}{2}) \\
\text{ambient\_diffuse} &= (1{-}F_{\text{env}})(1{-}m)\,\text{albedo}\,
  \big(\text{ssgi\_strength}\cdot\text{ssgi} +
  \text{env\_diffuse}\cdot\text{ao\_factor}\big) \\
\text{ambient\_specular} &= F_{\text{env}}\cdot
  \operatorname{mix}(\text{env\_specular},\ \text{ssr},\
  \text{smoothness}^2)
\end{align*}
and both terms are added to \code{lit}. The \code{ambient} slider is
vestigial as a shading input: the CPU folds it into
\code{sky\_color}/\code{ground\_color} as a uniform grey lift before
upload (\code{render.rs}), so the environment hemisphere is the single
source of non-bounce ambient light.

Pass order per frame: Hi-Z $\to$ GTAO (+temporal, blur) $\to$ SSGI
(+temporal, blur) $\to$ SSR (+temporal) $\to$ deferred shade $\to$ blit.

\section{Outputs, tone mapping, debug modes}

The pass writes \textbf{two} targets: \code{display} (what the user sees,
subject to debug overrides) and \code{lit} (the \emph{linear, pre-tonemap,
pre-exposure} shaded color, clamped at zero and rotated into a history
texture that feeds next frame's SSGI --- debug modes never pollute the
feedback loop, and the bounce integrates HDR radiance rather than
tone-compressed values). Tone mapping is the ACES Narkowicz fit applied to
$\text{lit} \times \text{exposure}$ inline in the shader for
\code{display}; gamma is deferred to \code{blit.wgsl}'s exact piecewise
sRGB curve.

Nineteen debug modes (\code{DBG_*}) can override \code{display}: albedo
(raw/true), normals, roughness, packed PBR, depth, specular-only,
diffuse-only, average shadow transmittance, cubemap-face color, plaster
(albedo replaced by white), alpha, primitives-only, DSM moments, AO, SSGI,
lit history, SSR, and depth standard deviation
$\sigma = \sqrt{E[z^2] - E[z]^2}$ from the moment channels.

\section{The non-deferred path: the compositor}

When deferred shading is off, volumes and primitives render to separate
color targets and \code{rmesh-compositor}'s \code{composite.wgsl} merges
them per pixel: recover the volume's view depth by un-premultiplying its
expected-depth channel, recover the primitive depth from the hardware
buffer, then either pass the volume through (no primitive), output the
primitive opaquely (primitive in front), or composite the premultiplied
volume \emph{over} the primitive:
$c = c_{\text{tet}} + (1 - \alpha_{\text{tet}})\, c_{\text{prim}}$.
This is a single-layer depth resolve --- correct because the volume was
already composited in sorted order upstream and primitives are opaque.

The opaque primitives themselves (cube, sphere, plane, cylinder, custom glTF
meshes) are drawn by \code{PrimitivePipeline}
(\code{primitive_pass.rs}) with full glTF metallic-roughness materials
(base-color, metal-rough, normal, and occlusion textures via a
\code{MaterialRegistry}), in either single-target mode or 4-target MRT mode
matching the G-buffer above.

